% !TEX root = ../main.tex
\corrigesChapitre{5}{Fonction exponentielle}
\corrigesSouples

\begin{blocLivre}{cor-5-1}
\begin{corrige}{1}{Simplifier une expression}
\demarche{Dans un produit, additionner les exposants ; dans un quotient, les soustraire.}
On applique les règles de calcul avant de simplifier les exposants :
\begin{align*}
e^{2x}e^{3-x}&=e^{2x+3-x}=\res{e^{x+3}},\\
\frac{e^{5x}}{e^{2x-1}}&=e^{5x-(2x-1)}=\res{e^{3x+1}},\\
(e^{-x})^3e^{2x}&=e^{-3x}e^{2x}=\res{e^{-x}}.
\end{align*}
Dans le quotient, les parenthèses conservent le signe du terme $-1$.

\end{corrige}
\end{blocLivre}
\begin{blocLivre}{cor-5-2}
\begin{corrige}{2}{Des équations immédiates}
\demarche{Utiliser l’injectivité de l’exponentielle et sa stricte positivité.}
L'exponentielle est strictement croissante, donc injective :
$e^{2x-1}=e^3\iff2x-1=3\iff\res{x=2}$.
Puis $e^{x+2}=1=e^0\iff\res{x=-2}$.
Enfin $e^x>0$ pour tout $x$, donc $e^x=-2$ n'a $\res{\text{aucune solution réelle}}$.
\end{corrige}
\end{blocLivre}
\begin{blocLivre}{cor-5-3}
\begin{corrige}{3}{Comparer des exponentielles}
\demarche{Comparer les exposants grâce à la croissance stricte.}
Par croissance stricte, $e^{3x+1}\leq e^{x+5}\iff3x+1\leq x+5\iff\res{x\leq2}$.
$e^{-2x}>1=e^0\iff-2x>0\iff\res{x<0}$.
Comme $-1<0<1/2$, on a $\res{e^{-1}<1<e^{1/2}}$.
\end{corrige}
\end{blocLivre}
\begin{blocLivre}{cor-5-4}
\begin{corrige}{4}{Dériver et donner le signe}
\demarche{Le signe de la constante devant l’exponentielle détermine celui de la dérivée.}
$f'(x)=2e^{2x}>0$ : $f$ est strictement croissante.
$g'(x)=-3e^{-x}<0$ : $g$ est strictement décroissante.
$h'(x)=e^{x+1}>0$ : $h$ est strictement croissante.
Ces résultats valent sur tout $\R$.
\end{corrige}
\end{blocLivre}
\begin{blocLivre}{cor-5-5}
\begin{corrige}{5}{La tangente en zéro}
\demarche{Une tangente fournit une approximation locale, avec un signe d’approximation.}
$f(0)=1$ et $f'(0)=1$, donc $\res{T:y=x+1}$.
Pour un petit $h$, $e^h\approx1+h$ : $\res{e^{0{,}03}\approx1{,}03}$.
Il s'agit d'une approximation, pas d'une égalité.
\end{corrige}
\end{blocLivre}
\begin{blocLivre}{cor-5-6}
\begin{corrige}{6}{Factoriser avant de résoudre}
\demarche{Factoriser l’exponentielle, toujours strictement positive.}
\begin{questions}
\item $e^{2x}-e^x=e^x(e^x-1)$. Le premier facteur ne s'annule pas :
$e^x=1\iff\res{x=0}$.
\item $(x-2)e^x=0\iff\res{x=2}$.
\item $e^{-x}>0$, donc le signe est celui de
$x^2-3x+2=(x-1)(x-2)$. L'ensemble des solutions est $\res{[1;2]}$.
\end{questions}
\end{corrige}
\end{blocLivre}
\begin{blocLivre}{cor-5-7}
\begin{corrige}{7}{Étudier un produit}
\demarche{Factoriser la dérivée par l’exponentielle et comparer les valeurs aux bornes.}
\begin{questions}
\item Par produit, $f'(x)=e^{-x}-(x+1)e^{-x}=\res{-xe^{-x}}$.
\item Le signe de $f'$ est celui de $-x$.
\end{questions}
\begin{center}\begin{tableau}
\otTab[3.1]{$x$/0.8,$f'(x)$/0.8,$f$/1.8}{$-2$,$0$,$4$}
\tkzTabLine{,+,z,-,}
\tkzTabVar{-/$-e^2$,+/$1$,-/$5e^{-4}$}
\end{tableau}\end{center}
Le maximum est $\res{1}$ en 0. Le minimum sur le segment vaut $\res{-e^2}$ en $-2$,
car l'autre borne donne une valeur positive.
\par\otLabel{3.}\enspace $f(1)=2/e$ et $f'(1)=-1/e$, donc
$\res{T:y=(3-x)/e}$. Le signe de $f$ est celui de $x+1$ : sur le domaine,
$\res{f(x)\geq0\iff x\in[-1;4]}$.
\end{corrige}
\end{blocLivre}
\begin{blocLivre}{cor-5-8}
\begin{corrige}{8}{Une différence d'exponentielles}
\demarche{Faire apparaître soit un produit, soit un carré positif.}
\begin{questions}
\item $g'(x)=2e^{2x}-2e^x=2e^x(e^x-1)$.
Elle est négative pour $x<0$, nulle en 0, positive pour $x>0$.
\item $g$ décroît puis croît : son minimum vaut $\res{g(0)=-1}$.
\item $(e^x-1)^2\geq0$, donc $g(x)\geq-1$ avec égalité exactement lorsque $e^x=1$,
c'est-à-dire $x=0$.
\end{questions}
\label{figlivre-r05-08-var-cor}
\figureLivre[corrige]{r05-08-var}

\end{corrige}
\end{blocLivre}
\begin{blocLivre}{cor-5-9}
\begin{corrige}{9}{Un quotient toujours défini}
\demarche{La positivité du dénominateur autorise les calculs et la comparaison des fractions.}
\begin{questions}
\item Le dénominateur $1+e^x$ est strictement positif sur $\R$.
Les fonctions en jeu sont dérivables, donc le quotient aussi.
\item $f'(x)=\dfrac{e^x(1+e^x)-e^xe^x}{(1+e^x)^2}
=\res{\dfrac{e^x}{(1+e^x)^2}}>0$ : croissance stricte sur $\R$.
\item Le numérateur est positif et inférieur au dénominateur, donc $0<f<1$.
$f(x)=1/2\iff2e^x=1+e^x\iff e^x=1\iff\res{x=0}$.
\end{questions}
\end{corrige}
\end{blocLivre}
\begin{blocLivre}{cor-5-10}
\begin{corrige}{10}{Passer du continu au discret}
\demarche{Échantillonner le modèle continu aux instants entiers.}
\begin{questions}
\item $N(0)=\res{500}$ et $N(5)=\res{500e\approx1359}$ individus.
L'arrondi concerne la valeur calculée ; le modèle continu peut donner un nombre non entier.
\item $u_n=500(e^{0{,}2})^n$ : suite géométrique de premier terme 500 et
de raison $\res{e^{0{,}2}}$.
\item Le taux annuel vaut $e^{0{,}2}-1\approx0{,}2214028$, soit
$\res{22{,}1\,\%}$ à 0,1 point près. Ce taux n'est pas 20\,\%.
\end{questions}
\end{corrige}
\end{blocLivre}
\begin{blocLivre}{cor-5-11}
\begin{corrige}{11}{Construire une approximation d'Euler}
\demarche{Chaque étape d’Euler utilise la pente au début de cette étape.}
\begin{questions}
\item $y_{k+1}=(1+1/N)y_k$, donc la suite est géométrique de raison $1+1/N$.
Après $N$ étapes, $\res{y_N=(1+1/N)^N}$.
\item Pour $N=2$ : $y_2=(3/2)^2=2{,}25$ ; pour $N=10$ :
$y_{10}=1{,}1^{10}\approx2{,}593742$.
Les deux valeurs sont inférieures à $e\approx2{,}718282$ ; celle avec dix étapes est plus proche.
Les termes de la suite $\big((1+1/n)^n\big)$ fournissent ainsi des valeurs approchées de $e$.
\item Une boucle reproduit directement l'actualisation d'Euler :
\end{questions}
\begin{python}
def euler(N):
    h = 1/N
    y = 1.0
    for k in range(N):
        y = y + h*y
    return y
\end{python}

Les $N$ étapes de longueur $1/N$ couvrent l’intervalle de 0 à 1. Dans la mise à jour, la pente utilisée est la valeur courante de $y$, puisque la fonction étudiée vérifie $y^\prime=y$.

\label{figlivre-r05-11-euler-cor}
\figureLivre[corrige]{r05-11-euler}

\end{corrige}
\end{blocLivre}
\begin{blocLivre}{cor-5-12}
\begin{corrige}{12}{Le refroidissement d'une boisson}
\demarche{Justifier la décroissance avant de chercher la première minute admissible.}
\begin{questions}
\item $T(0)=90\,{}^\circ$C.
$T'(t)=-7e^{-0{,}1t}<0$ : la température est strictement décroissante.
\item $70e^{-0{,}1t}>0$, donc $T(t)>20$ pour tout instant fini $t\geq0$.
\item $T(12)\approx41{,}08\,{}^\circ$C et $T(13)\approx39{,}08\,{}^\circ$C.
La première minute entière où $T\leq40$ est donc $\res{13}$.
\item Le programme suivant teste les minutes entières dans l'ordre :
\end{questions}
\begin{python}
from math import exp
t = 0
while 20 + 70*exp(-0.1*t) > 40:
    t = t + 1
\end{python}
À la sortie de la boucle, \code{t} vaut 13.
\end{corrige}
\end{blocLivre}
\begin{blocLivre}{cor-5-13}
\begin{corrige}{13}{Une concentration après injection}
\demarche{La dérivée permet de localiser et de prouver le pic de concentration.}
\begin{questions}
\item $C'(t)=8e^{-0{,}5t}-4te^{-0{,}5t}=\res{4(2-t)e^{-0{,}5t}}$.
La concentration croît sur $[0;2]$, puis décroît sur $[2;8]$.
Les valeurs remarquables sont $C(0)=0$, $C(2)=16/e$, $C(8)=64e^{-4}$.
\item Maximum à $\res{t=2\text{ h}}$, de valeur $\res{16/e\approx5{,}886\text{ mg/L}}$.
\item $C(0)=0$ et $C'(0)=8$ : la tangente est $\res{y=8t}$.
La pente est une vitesse initiale d'augmentation de 8 mg/L par heure.
\item Le taux moyen vaut $(C(3)-C(2))/(3-2)=24e^{-1{,}5}-16/e$
$\approx\res{-0{,}53\text{ mg/L par heure}}$. Son signe négatif traduit une diminution.
\end{questions}
\label{figlivre-r05-13-var-cor}
\figureLivre[corrige]{r05-13-var}

\label{figlivre-r05-13-concentration-cor}
\figureLivre[corrige]{r05-13-concentration}

\end{corrige}
\end{blocLivre}
\begin{blocLivre}{cor-5-14}
\begin{corrige}{14}{Prouver une inégalité sans convexité}
\demarche{Étudier la différence avec la droite pour prouver une inégalité sur tout le domaine.}
\begin{questions}
\item $F'(x)=e^x-1$, négative pour $x<0$, nulle en 0 et positive pour $x>0$.
\item $F$ décroît jusqu'à 0 puis croît : son minimum est $F(0)=0$.
Donc $F(x)\geq0$, soit $\res{e^x\geq1+x}$.
La stricte monotonie de part et d'autre assure que l'égalité a lieu seulement pour $x=0$.
\item Pour $x=0{,}1$, on obtient $e^{0{,}1}\geq1{,}1$ ; pour $x=-0{,}1$,
on obtient $e^{-0{,}1}\geq0{,}9$.
\item Comme les nombres sont positifs, $1/e^{0{,}1}\geq0{,}9$ entraîne
$e^{0{,}1}\leq1/0{,}9=10/9$. Ainsi $\res{1{,}1\leq e^{0{,}1}\leq10/9}$.
\item L'approximation affine en 0 vaut 1,1. Elle est ici une borne inférieure.
L'encadrement garantit une erreur d'au plus $10/9-1{,}1=1/90\approx0{,}0111$.
\end{questions}
\label{figlivre-r05-14-var-cor}
\figureLivre[corrige]{r05-14-var}

\label{figlivre-r05-14-inegalite-cor}
\figureLivre[corrige]{r05-14-inegalite}

\end{corrige}
\end{blocLivre}
\begin{blocLivre}{cor-5-15}
\begin{corrige}{15}{Choisir une durée de campagne}
\demarche{Étudier les variations, puis convertir les centaines d’euros en euros.}
\begin{questions}
\item Par produit, $G'(t)=100e^{-0{,}5t}-50(t+1)e^{-0{,}5t}
=\res{50(1-t)e^{-0{,}5t}}$.
\item La dérivée a le signe de $1-t$ : croissance sur $[0;1]$, décroissance sur $[1;6]$.
Le maximum est atteint après $\res{1\text{ semaine}}$ et vaut
$200e^{-1/2}$ centaines d'euros, soit $\res{20000/\sqrt e\approx12130{,}61\text{ euros}}$.
\item $G(0)=100$ et $G(1)=200/\sqrt e>100$ car $e<4$.
\item $G(2)=300/e\approx110{,}36$ et $G(3)=400e^{-1{,}5}\approx89{,}25$.
Sur $]0;1]$, la croissance stricte depuis $G(0)=100$ donne $G>100$.
Sur $[1;2]$, la décroissance donne $G(t)\geq G(2)=300/e>100$ car $e<3$.
\item Les durées entières favorables sont $\res{1\text{ et }2\text{ semaines}}$.
À partir de 3, le gain est déjà inférieur à 100 et continue de décroître jusqu'à 6.
\end{questions}
\label{figlivre-r05-15-var-cor}
\figureLivre[corrige]{r05-15-var}

\label{figlivre-r05-15-campagne-cor}
\figureLivre[corrige]{r05-15-campagne}

\end{corrige}
\end{blocLivre}
